\documentclass[../main.tex]{subfiles}
\begin{document}
%==========================================TIÊU ĐỀ TẬP========================================
\begin{table}[h]
  \begin{adjustwidth}{-1cm}{}
    \begin{tabular}{>{\centering\arraybackslash}p{6cm}|>{\raggedright\arraybackslash}p{12.5cm}}
      \multirow{5}{6cm}
      % Cột bên trái: ÉPISODE và số tập
      {\\[-40pt]
        \flaregothic\fontsize{20pt}{20pt}\selectfont \centering
        \color{eptype}ÉPISODE\color{black}\\
        \burbank\fontsize{72pt}{72pt}\selectfont
        \color{epnum}21\color{black}\\[6pt]
        \cafeteria\fontsize{20pt}{20pt}\selectfont\color{epnum}(2-4)\color{black}
        \\[18pt]
        \flaregothic\fontsize{14pt}{14pt}\selectfont
        \color{epattr}BIENVENUE, \uline{\color{Yubiyubi} Korone} !
        \flaregothic\fontsize{18pt}{18pt}\selectfont
        \color{epattr}ÉPISODE PERDU
        \color{black}
      }
       &
      % Dòng thứ nhất: tên tập tiếng Nhật
      {\fotskip\fontsize{18pt}{18pt}\selectfont \ruby{犬}{いぬ}\ruby{三}{さん}\ruby{銃}{じゅう}\ruby{士}{し}を\ruby{連}{つ}れてきたよ
      } \\[6pt]
      % Dòng thứ hai: tên tập tiếng Anh
       & {\cafeteria\fontsize{18pt}{18pt}\selectfont The Three Musketeers of Dogs are Here!}                     \\[4pt]
      % Dòng thứ ba: tên tập tiếng Việt
       & {\vnmsans\fontsize{18pt}{18pt}\selectfont Ba nàng khuyển ngự lâm}                                       \\[8pt]
      % Dòng thứ tư: Nhân vật xuất hiện
       & {\caviar{\fontsize{14pt}{14pt}\selectfont Nhân vật: \emp{kuon1} \emp{ryuuhou1} \emp{mio} \emp{korone}}} \\[8pt]
      % Dòng cuối cùng: Thời gian diễn ra của tập (thường sẽ là ngày phát hành)
       & {\continuum\fontsize{16pt}{16pt}\selectfont Ngày phát hành: 29/09/2019}                                 \\[18pt]
    \end{tabular}
  \end{adjustwidth}
\end{table}
%=========================================TRUYỆN CHÍNH========================================
\color{black}

\txtbx[2]{{\color{vol} \uline{Tóm tắt:}} Korone trổ tài câu cá ngay giữa lòng thành phố và thu về vô số món đồ kỳ quặc, bao gồm cả một thực thể ``Ác tính của con người''. Sự việc càng thêm rắc rối khi một chú chó hoang nuốt phải quả cầu bí ẩn và biến thành người chó, dẫn đến màn bắt nhầm đầy trớ trêu của sở thú dành cho nàng Khuyển nhà hololive.}

\pov{Hikawa Kuon}

\model{23}{r}{7cm}{korone}

Một buổi sáng cuối tháng Chín, không khí se lạnh nhưng ngập tràn nắng vàng. Tôi đang ngồi trên chiếc ghế dựa trong văn phòng, lười biếng nhìn ra phía cửa sổ. Ở nơi này, một ngày mới bắt đầu đồng nghĩa với việc bạn phải chuẩn bị sẵn tinh thần cho một điều gì đó điên rồ sắp xảy ra.

Và quả thực, linh cảm của tôi chẳng sai bao giờ. Bên ngoài cửa sổ, một cảnh tượng mà bộ não bình thường của tôi không tài nào tiêu hóa nổi đang diễn ra: một cô nàng với đôi tai chó màu nâu đang điềm nhiên ngồi trên bậu cửa sổ tầng trệt, tay lăm lăm một chiếc cần câu cá.

\chrint

Đó là Inugami Korone (\ruby{戌}{いぬ}\ruby{神}{がみ}ころね, \ruby{I-nư}{Tuất}-\ruby{ga-mi}{Thần}\ Cô-rô-ne), thành viên của \hll\ GAMERS. Một cô nàng Khuyển thực thụ - đúng như cái tên của mình - vốn là người quảng cáo cho một tiệm bánh mì ở Tokyo. Nghe thì có vẻ là một công việc khá bình lặng, nhưng nó lại cực kỳ hòa hợp với tính cách năng động và có chút kỳ quặc của cô. Korone cũng là một người cực kỳ yêu thích việc ngắm nhìn những chú chó đồng loại với mình.

Cô mang đầy đủ những đặc tính đáng yêu của một chú chó: ngọt ngào, đôi khi hơi ngốc nghếch, và đặc biệt là vô cùng trung thành với những người mình tin tưởng, tiêu biểu là cô bạn thân Okayu và cả tôi nữa (người mà cô ấy nhất quyết coi là chủ nhân, dù tôi đã bảo là không cần thiết phải gọi như thế).

Korone có niềm đam mê mãnh liệt với các trò chơi cổ điển từ thập niên 80, 90, nhất là những tựa game ít người biết đến hoặc thậm chí là những trò chơi kém chất lượng (\textit{kusoge}) và kỳ lạ thay, cô lại tìm thấy niềm vui bất tận từ chúng.

Tuy nhiên, đừng để vẻ ngoài dễ thương đó đánh lừa. Đôi lúc, nàng Khuyển này có thể trở nên đáng sợ một cách vô tư đến mức khiến người ta lạnh sống lưng. Cô có thể cười nói vui vẻ khi sử dụng cưa máy trong game kinh dị hay chẳng hề nao núng trước những cảnh rùng rợn. Chính sự hồn nhiên pha lẫn chút ``điên rồ'' đó đã tạo nên một sức hút khó cưỡng cho Korone trong lòng người hâm mộ.

Về ngoại hình, Korone sở hữu mái tóc màu nâu nhạt dài vừa phải với hai bím tóc tết lại, cùng đôi tai chó cụp xuống trông như giống chó Spaniel. Cô diện một chiếc áo khoác vàng khoác trễ vai bên ngoài chiếc váy trắng, đôi mắt nâu tông-sơ-tông với màu tóc luôn tỏa sáng vẻ tinh nghịch. Chiếc vòng cổ cho chó màu đỏ chính là điểm nhấn hoàn hảo cho vẻ ngoài đặc trưng của cô.

\chrint

\img{2-5a}

Bụng của nàng Chó đang kêu ọc ọc vì đói, thế nhưng em ấy vẫn kiên nhẫn cầm cần câu, đôi chân đung đưa trong không khí đầy thong thả. ``Hem biết có câu được gì nữa không ta\~{}'' Korone ngân nga, vẻ mặt hồn nhiên như một đứa trẻ đang đi dã ngoại.

Ngồi cạnh tôi lúc này là Ryuuhou. Con bé đang chống cằm nhìn chằm chằm ra cửa sổ, đôi mắt hai màu lấp lánh sự tò mò. Nó huých tay tôi, thì thầm: ``Này Onii-san, chị ấy ngồi đó nửa tiếng rồi đấy. Anh nghĩ chị ấy định câu cá mập trên vỉa hè à?''

Tôi bĩu môi, cố nhịn cười: ``Anh thấy từ nãy đến giờ toàn câu được mấy thứ gì đâu không à. Dưới đường nhựa làm gì có cá mà câu cơ chứ.''

Đúng lúc đó, nàng Sói đen Mio xuất hiện từ phía sau, khiến cả tôi và Ryuuhou đều giật mình. Mio nghiêng đầu nhìn xuống Korone, rồi quay sang tôi với vẻ mặt đầy hoang mang: ``Đúng rồi, chúng ta đang ở giữa thành phố mà, câu bằng niềm tin à? Mà Kuon-senpai bảo `mấy thứ gì đâu' là sao?''

Tôi không đáp lời, chỉ lặng lẽ chỉ tay về phía người đang cầm cần câu, ra dấu cho Mio và Ryuuhou hãy tự mình chứng kiến.

Ngay lập tức, Korone liên tiếp kéo cần lên với những cú giật mạnh mẽ. Cả ba chúng tôi đều đứng hình khi thấy những thứ lần lượt xuất hiện ở đầu dây câu: từ một tấm áp phích phim cũ nát, một chiếc laptop đời cổ, một cốc trà sữa uống dở cho đến một con thú bông nhỏ xíu. Đỉnh điểm là một chú chó hoang trưởng thành cũng bị em ấy kéo lên một cách thần kỳ!

\gif{2-5b}{1}{107}

Ryuuhou phấn khích reo lên: ``Oa! Đỉnh thật sự! Chị ấy vừa câu được một con chó kìa!''

Korone quay lại cười tươi rói, đuôi vẫy rối rít: ``Cơ mà bà câu được nhiều thứ lạ lắm nha! Không chỉ đồ vật mà còn tìm được cả bạn mới nữa nè!''

Mio thì há hốc mồm, giọng đầy kinh ngạc: ``Thế quái nào mà bà lại câu được mấy thứ đó hả Korone!?''

Tôi chỉ biết bất lực nhún vai, đáp lại cái nhìn cầu cứu của Mio: ``Đừng hỏi anh. Anh cũng biết chết liền á. Logic của văn phòng này vốn đã tuyệt chủng từ lâu rồi.''

Bất chợt, Korone lại reo lên một cách phấn khích. Chiếc cần câu trong tay em ấy đột ngột cong vút xuống, rung lắc dữ dội như thể vừa mắc phải một con quái vật khổng lồ dưới lòng đất.

\img{2-5c1}

``\textbf{Ồ, em câu được một con bự nè!}'' Korone hét lớn, gồng mình kéo cần. 

Ryuuhou cũng hùa theo cổ vũ nhiệt tình: ``Cố lên chị ơi! Kéo mạnh lên!'' 

Trái ngược với sự hào hứng của hai người họ, Mio lại khoanh tay, nhíu mày đầy lo lắng: ``Tôi cảm thấy có gì đó không ổn chút nào đâu nha\dots''

\img{2-5c2}

Chưa kịp dứt lời, Korone đột nhiên kêu lên một tiếng ``Ơ\dots!'' đầy ngỡ ngàng. Sức kéo từ bên dưới quá lớn khiến em ấy bị nhấc bổng lên khỏi bậu cửa sổ, lơ lửng giữa không trung một giây rồi bị kéo tuột xuống dưới, biến mất khỏi tầm mắt chúng tôi.

Cả ba chúng tôi đứng chết lặng. Tôi thảng thốt kêu lên: ``Vãi thật! Cá voi cắn câu hay gì?''

Mio hốt hoảng túm lấy áo tôi: ``Em không biết nữa, nhưng phải đi cứu em ấy ngay thôi!'' 

Ryuuhou thì đã nhanh chân chạy ra phía cửa: ``Đi thôi Onii-san, Mio-san! Mau lên kẻo chị ấy bị quái vật vỉa hè ăn thịt mất!''

Thế nhưng, khi chúng tôi vừa định lao ra ngoài thì cánh cửa phòng lại bật mở. Korone xuất hiện trở lại, tươi cười rạng rỡ như chưa từng có cuộc chia ly nào, trên tay vẫn cầm chiếc cần câu báu vật. ``Tôi về rồi đây\~{}''

\img{2-5d}

Mio đứng hình mất vài giây rồi lắp bắp hỏi: ``Quái, bà vẫn ổn à Korone!? Sao lên nhanh thế?''

Korone điềm nhiên giải thích: ``À, lúc rơi xuống bà bám vào được cái ống nước bên cạnh nên leo lên đây luôn. Bà vẫn ổn mà, Kuon-senpai, Ryuuhou-chan!''

Tầm mắt chúng tôi lập tức đổ dồn vào thứ mà Korone vừa mang về, đang lủng lẳng ở đầu dây câu. Đó là một quả cầu màu tím, bề mặt nhẵn bóng và phát ra thứ ánh sáng kỳ quái.

Mio nheo mắt nhìn kỹ: ``Mà nhân tiện, bà nhặt được cái gì thế này?''

Nàng Khuyển ngây ngô đáp: ``Tôi cũng hổng biết nữa. Nó cứ tự dưng nổi lên khi tôi kéo dây câu thôi à.''

\img{2-5e}

Tôi tiến lại gần, rùng mình khi đọc được dòng chữ màu trắng kỳ lạ trên bề mặt quả cầu: ``Ác tính của con người''.

Ryuuhou tò mò định đưa tay chạm vào nhưng tôi kịp thời ngăn lại. Em tôi lẩm bẩm: ``Ác tính của con người? Nghe giống tên mấy cái vật phẩm trong game RPG em hay chơi thế.''

Mio thì tỏ vẻ lo ngại thật sự: ``Cái tên nghe chẳng lành chút nào\dots''

Bất thình lình, chú chó hoang mà Korone câu được lúc nãy nhảy phốc tới, sủa ầm ĩ rồi ngoạm lấy quả cầu tím, chạy tọt vào góc phòng. Trước sự ngỡ ngàng của cả bọn, nó gặm quả cầu điên cuồng rồi nuốt chửng luôn vào bụng.

Tôi đổ mồ hôi hột: ``Nó vừa ăn cái đó thật luôn à?'' 

Mio thở dài, giọng đầy vẻ tiên tri: ``Vâng\dots\ em nghĩ chuyện này sẽ chẳng dẫn tới điều gì hay ho đâu.''

Quả nhiên, chú chó hoang bắt đầu phát sáng rực rỡ. Một luồng ánh sáng tím bao phủ lấy nó, khiến chúng tôi phải lấy tay che mắt. Khi ánh sáng dịu đi, đứng trước mặt chúng tôi không còn là một chú chó bình thường, mà là một sinh vật nửa người nửa chó, đứng thẳng trên hai chân và mặc một chiếc quần đen.

\img{2-5f}

Nó cất tiếng nói bằng một giọng đều đều như robot: ``Vì tôi đã ăn `Ác tính của con người' nên giờ tôi biến thành người chó rồi nhe\dots''

Ryuuhou há hốc mồm, chỉ tay vào sinh vật đó: ``Oa! Tiến hóa kìa Onii-san! Anh nhìn kìa!''

Tôi thì chỉ biết bóp trán: ``Cái logic ảo ma gì thế này trời\dots''

Trong khi chúng tôi còn đang sốc, Korone đã nhảy bổ lên phía trước. Em ấy chống nạnh, nhìn chằm chằm vào kẻ lạ mặt và dõng dạc tuyên bố: ``Tôi đây cũng là người chó nè! Chỗ này là địa bàn của tôi, ông biến ra ngoài đi!''

Người chó giật mình, vội vàng cúi gập người xin lỗi nàng Khuyển một cách rất lịch sự. Mio đứng đơ như tượng, còn tôi chỉ biết thở dài tự nhủ: ``\textit{Thì cả hai đều là người chó thật mà, nhưng sao cảnh này nó cứ sai sai thế nào ấy\dots}''

Không để tình hình kéo dài thêm, Korone rút điện thoại ra, bấm số gọi với vẻ mặt cực kỳ nghiêm túc: ``A-lô, sở bắt thú đấy ạ? Ở đây có một con người chó màu nâu đi lạc, phiền các anh đến hốt đi dùm nhe!''

Tôi giật nảy mình, vội vàng ngăn lại: ``Này em ơi, em nói thế người ta sẽ hiểu lầm to cho xem!''

\ct

Một lát sau, nhân viên sở thú có mặt. Đó là một người phụ nữ trung niên với bộ đồng phục chỉnh tề và ánh mắt sắc lẹm. Korone hớn hở chạy ra vẫy tay: ``Chào chị, em gọi đây ạ!''

Nhưng không nói lời nào, người phụ nữ kia lập tức tiến tới, còng tay Korone lại bằng một động tác chuyên nghiệp đến mức đáng sợ. 

\img{2-5g}

Cả đám chúng tôi, kể cả tên người chó, đều sững sờ. Korone ngơ ngác kêu lên một tiếng ``Ể?'', nhưng đã quá muộn để phản kháng. Em ấy bị lôi đi trong sự ngỡ ngàng tột độ của tất cả mọi người.

Tôi đứng im tại chỗ, cảm thấy não mình vừa bị quá tải nghiêm trọng. Tôi bật ra một tiếng cười đau khổ, quay sang Mio và Ryuuhou: ``Thấy chưa? Anh đã bảo mà! Người ta bắt `người chó' là bắt em ấy chứ bắt cái tên kia làm gì!''

Ryuuhou gãi đầu, nhìn theo chiếc xe sở thú vừa lăn bánh: ``Tội nghiệp Korone-san, nhưng mà\dots\ nhìn mặt chị ấy lúc bắt đi trông buồn cười thật.''

Nàng Sói đen lắc đầu ngán ngẩm, trong khi tên người chó vẫn đứng đó, ngơ ngác không hiểu tại sao mình lại là kẻ thoát nạn.

%==========================================TRUYỆN PHỤ=========================================
\omk

Sau khi Korone bị lôi đi, ba chúng tôi đứng nhìn nhau một lúc lâu giữa căn phòng tĩnh lặng đến lạ thường. Cuối cùng, Mio thở dài một hơi thườn thượt, lay lay vai tôi để kéo tôi ra khỏi trạng thái hóa đá.

Cô nàng tai thú lo lắng hỏi: ``Giờ phải làm sao đây anh? Chúng ta không thể để Korone ở trong sở thú được.''

Ryuuhou thì lại có vẻ hào hứng lạ thường: ``Đi giải cứu chị ấy thôi Onii-san! Em muốn xem chị ấy ở trong chuồng trông sẽ như thế nào!''

Tôi gõ nhẹ vào đầu con bé một cái: ``Nghịch ngợm vừa thôi. Anh phải tìm cách giải thích với bên đó, chứ để lâu là em ấy thành `vật mẫu triển lãm' của người ta luôn bây giờ.''

\ct

Chiều hôm đó, tôi, Mio và cả tên người chó kỳ lạ kia (hắn nhất quyết đòi đi theo để tạ lỗi) kéo nhau đến sở thú. Ryuuhou cũng nằng nặc đòi đi bằng được, miệng liên tục lẩm bẩm về một ``chiến dịch giải cứu anh hùng''.

Vừa bước vào khu vực tiếp khách, cảnh tượng đầu tiên đập vào mắt khiến tôi suýt nữa ngã ngửa. Nàng Khuyển Korone đang ngồi thong thả trên ghế, hai tay cầm máy Switch chơi game hăng say, thi thoảng còn reo lên đầy phấn khích. Chung quanh em ấy là vài nhân viên sở thú đang đứng nhìn với vẻ mặt vừa tò mò vừa thán phục.

\img{2-5h}

Ryuuhou chạy tới, tròn mắt nhìn: ``Korone-san! Chị đang chơi game trong sở thú thật đấy à?''

Korone ngẩng đầu lên, cười tươi rói: ``À, Ryuuhou-chan! Ở đây phục vụ tốt lắm nha, có cả wifi mạnh nữa nè!''

Mio thở phào nhẹ nhõm, ngồi xuống cạnh em ấy rồi hai người bắt đầu buôn chuyện đủ thứ trên đời, từ đồ ăn sáng cho đến các trò chơi cổ điển. Nhìn họ, tôi cứ ngỡ mình đang ở quán cà phê chứ không phải nơi giam giữ thú hoang.

Trong khi đó, tôi phải vất vả đi tìm quản lý sở thú để giải thích về sự nhầm lẫn tai hại này. Việc chứng minh Korone là một thần tượng \hll\ chứ không phải sinh vật biến dị thực sự là một thử thách khó khăn, nhất là khi tên người chó mặc quần đen cứ lù lù đứng bên cạnh tôi làm bằng chứng sống cho sự kỳ quặc.

Cuối cùng, sau một hồi thuyết phục (và cả một khoản phí bồi thường không nhỏ khiến ví tiền tôi khóc thét), Korone cũng được thả. Bước ra khỏi sở thú, em ấy ngúng nguẩy đuôi, hớn hở nói: ``Ở trong đó trải nghiệm mới lạ thật đấy, cảm ơn mọi người đã đến đón em nha!''

Ryuuhou khoác tay Korone, cười tít mắt: ``Lần sau chị lại vào đó thì nhớ gọi em đi cùng với nhe!''

Tôi chỉ biết lắc đầu ngao ngán, nhìn ba bóng dáng - một nàng Sói, một nàng Chó và một cô em gái tinh nghịch - đang đi phía trước. 

``\textit{Ba nàng khuyển ngự lâm này\dots\ đúng là bộ ba oái oăm nhất hành tinh mà mình từng gặp.}'' tôi thầm nghĩ, dù vậy, trên môi vẫn không nén nổi một nụ cười nhẹ. 

Có lẽ, những ngày tháng rắc rối này chính là thứ gia vị khiến cuộc sống của tôi tại \hll\ trở nên đáng nhớ hơn bao giờ hết.

\end{document}